\begin{thema}{Γ}

\begin{center}
\begin{tikzpicture}[scale=0.7]
  \draw[help lines, gray!30] (-3,-2) grid (5,4);
  \draw[-latex] (-3.3,0) -- (5.3,0) node[right] {$x$};
  \draw[-latex] (0,-2.3) -- (0,4.3) node[above] {$y$};
  \foreach \x in {-2,-1,1,2,3,4} \draw (\x,0.1) -- (\x,-0.1) node[below,font=\footnotesize] {$\x$};
  \foreach \y in {-1,1,2,3} \draw (0.1,\y) -- (-0.1,\y) node[left,font=\footnotesize] {$\y$};
  
  \draw[very thick, blue!70!black, smooth, tension=0.5]
    plot coordinates {(-2,0) (-1,1.5) (0,2) (0.5,2.7) (1,3) (2,1) (3,-1) (3.5,-0.5) (4,1)};
    
  \filldraw[blue!70!black] (-2,0) circle (2pt);
  \filldraw[blue!70!black] (1,3) circle (2pt) node[above,font=\footnotesize] {$(1,3) \ \text{Ολικό / Τοπικό Μέγιστο}$};
  \filldraw[blue!70!black] (3,-1) circle (2pt) node[below,font=\footnotesize] {$(3,-1) \ \text{Ολικό / Τοπικό Ελάχιστο}$};
  \filldraw[blue!70!black] (4,1) circle (2pt) node[right,font=\footnotesize] {$(4,1)$};
  
  % Chord and Tangent for MVT
  \draw[dashed, orange, very thick] (-2,0) -- (4,1);
  \node[orange!80!black, font=\footnotesize, fill=white, inner sep=1pt] at (2.5, 0.4) {$\text{Χορδή } \lambda=1/6$};
  
  \draw[dashed, red!60!black, very thick] (-1.5, 1) -- (2.5, 1.66) node[right, fill=white, inner sep=1pt] {$\text{Εφαπτομένη } \varepsilon \parallel \text{Χορδή}$};
  \filldraw[red!60!black] (0.5, 1.33) circle (2pt) node[below right, font=\footnotesize, fill=white, inner sep=1pt] {$\xi \approx 0.5$};
  
  % Bolzano y=1 line
  \draw[dashed, green!60!black, very thick] (-3,1) -- (5,1) node[right, fill=white, inner sep=1pt] {$y=1 \ (\text{3 τομές})$};
\end{tikzpicture}
\end{center}

\begin{erwthma}
\item Από την ανάλυση της γραφικής παράστασης της $f$ στο διάστημα $[-2, 4]$, παρατηρούμε ότι η συνάρτηση είναι:
\begin{itemize}
    \item $\textbf{γνησίως αύξουσα}$ στο διάστημα $[-2, 1]$.
    \item $\textbf{γνησίως φθίνουσα}$ στο διάστημα $[1, 3]$.
    \item $\textbf{γνησίως αύξουσα}$ στο διάστημα $[3, 4]$.
\end{itemize}
Συνεπώς, η $f$ παρουσιάζει:
\begin{itemize}
    \item Τοπικό μέγιστο στο $x=1$ με τιμή $f(1)=3$.
    \item Τοπικό ελάχιστο στο $x=3$ με τιμή $f(3)=-1$.
\end{itemize}
Για τον προσδιορισμό των ολικών ακροτάτων στο κλειστό $[-2, 4]$, συγκρίνουμε τις τιμές στα άκρα και στα τοπικά ακρότατα:
\[ f(-2)=0, \quad f(1)=3, \quad f(3)=-1, \quad f(4)=1 \]
Άρα:
\begin{itemize}
    \item Το $\textbf{ολικό μέγιστο}$ είναι το $3$ και επιτυγχάνεται για $x=1$.
    \item Το $\textbf{ολικό ελάχιστο}$ είναι το $-1$ και επιτυγχάνεται για $x=3$.
\end{itemize}

\item Η συνάρτηση $f$ είναι παραγωγίσιμη στο $[-2, 4]$, άρα είναι και συνεχής στο διάστημα αυτό. Εφαρμόζοντας το \eng{Lagrange} ΘΜΤ στο διάστημα $[-2, 4]$, υπάρχει $\xi \in (-2, 4)$ τέτοιο ώστε:
\[ f'(\xi) = \frac{f(4) - f(-2)}{4 - (-2)} \]
Αντικαθιστώντας τις τιμές $f(4)=1$ και $f(-2)=0$, έχουμε:
\[ f'(\xi) = \frac{1 - 0}{6} = \frac{1}{6} \]
Γραφικά, το $\xi$ αντιστοιχεί στο σημείο της $C_f$ όπου η εφαπτομένη είναι παράλληλη με την χορδή που ενώνει τα σημεία $(-2, 0)$ και $(4, 1)$. Από το σχήμα, αυτό συμβαίνει περίπου για $x \approx -0,5$ και $x \approx 3,5$.

\item Θεωρούμε τη συνάρτηση $g(x) = f(x) - 1$, η οποία είναι συνεχής στο $[-2, 4]$. Η εξίσωση $f(x)=1$ είναι ισοδύναμη με την $g(x)=0$.
Ελέγχουμε τα πρόσημα της $g$ σε διάσματα:
\begin{itemize}
    \item Στο διάστημα $[-2, 0]$: $g(-2) = 0 - 1 = -1 < 0$ και $g(0) = 2 - 1 = 1 > 0$. Από το \eng{Bolzano} θεωρήμα, υπάρχει τουλάχιστον μία λύση στο $(-2, 0)$.
    \item Στο διάστημα $[1, 3]$: $g(1) = 3 - 1 = 2 > 0$ και $g(3) = -1 - 1 = -2 < 0$. Από το \eng{Bolzano} θεωρήμα, υπάρχει τουλάχιστον μία λύση στο $(1, 3)$.
    \item Στο σημείο $x=4$: $g(4) = 1 - 1 = 0$, άρα το $x=4$ είναι λύση.
\end{itemize}
Επειδή η συνάρτηση είναι μονοτόνα σε αυτά τα διαστήματα, η εξίσωση $f(x)=1$ έχει ακριβώς $\textbf{3 λύσεις}$ στο $[-2, 4]$.

\item Το ορισμένο ολοκλήρωμα $\int_{-2}^{4} f(x) \d x$ αντιπροσωπεύει το πρόσημο της εμβαδικής περιοχής μεταξύ της $C_f$ και του άξονα $x'$.
Παρατηρούμε ότι:
\begin{itemize}
    \item Για $x \in [-2, 2,5]$ (περίπου), η $f(x) \ge 0$, οπότε το εμβαδόν είναι θετικό και αρκετά μεγάλο (με μέγιστο ύψος $3$).
    \item Για $x \in [2,5, 3,5]$ (περίπου), η $f(x) \le 0$, οπότε το εμβαδόν είναι αρνητικό, αλλά η περιοχή είναι μικρή (με ελάχιστο ύψος $-1$).
    \item Για $x \in [3,5, 4]$, η $f(x) \ge 0$, προσθέτοντας ξανά θετικό εμβαδόν.
\end{itemize}
Επειδή η περιοχή που βρίσκεται πάνω από τον άξονα $x'$ είναι εμφανώς μεγαλύτερη από την περιοχή που βρίσκεται κάτω από αυτόν, συμπεραίνουμε ότι:
\[ \int_{-2}^{4} f(x) \d x > 0 \]
\end{erwthma}
\end{thema}
%=========== ΤΕΛΟΣ ΛΥΣΗΣ - ΘΕΜΑ Γ ===========